%=============================================================================
\section{实战：强化学习实验矩阵}
%=============================================================================

现在我们把前面学到的知识串起来，实现一个完整的实验工作流。

我们的实验是一个三维矩阵：

\begin{itemize}[nosep]
  \item \textbf{维度 1}：多个 RL 环境（CartPole, LunarLander, HalfCheetah, \ldots）
  \item \textbf{维度 2}：多个算法/架构（PPO, SAC, DQN, \ldots）
  \item \textbf{维度 3}：每组跑 $N$ 个随机种子（确保统计显著性）
\end{itemize}

目标：一键提交所有实验，自动记录结果，方便后续分析。

\subsection{项目结构}

先来看整个项目长什么样：

\begin{lstlisting}
project/
  train.py              # 训练脚本，接受命令行参数
  configs/
    envs.txt            # 环境列表，每行一个
    archs.txt           # 架构列表，每行一个
  scripts/
    run_matrix.sh       # SLURM Array Job 脚本
    submit_all.sh       # 提交脚本
  results/              # 程序输出的结果
  logs/                 # SLURM 日志文件
\end{lstlisting}

\texttt{results/} 存放你程序写出的模型和指标，\texttt{logs/} 存放 SLURM 的标准输出和错误日志。两者分开管理比较清晰。

\subsection{配置文件}

用纯文本文件列出实验参数，方便修改，也方便脚本读取：

\begin{lstlisting}[title={\texttt{configs/envs.txt}}]
CartPole-v1
LunarLander-v2
HalfCheetah-v4
Hopper-v4
Walker2d-v4
\end{lstlisting}

\begin{lstlisting}[title={\texttt{configs/archs.txt}}]
PPO
SAC
DQN
TD3
A2C
\end{lstlisting}

\subsection{训练脚本接口}

\texttt{train.py} 需要接受命令行参数，这样 SLURM 脚本才能通过参数控制每次运行的内容：

\begin{lstlisting}[language=python,title={\texttt{train.py}（参数部分）}]
import argparse

parser = argparse.ArgumentParser()
parser.add_argument("--env", type=str, required=True)
parser.add_argument("--arch", type=str, required=True)
parser.add_argument("--seed", type=int, required=True)
parser.add_argument("--output-dir", type=str, default="results")
args = parser.parse_args()

# 结果保存到: results/{env}/{arch}/seed_{seed}/
\end{lstlisting}

关键设计：把环境、算法、种子都做成参数，这样同一个脚本就能覆盖所有实验组合。

%=============================================================================
\subsection{方案一：SLURM Array Job（推荐）}
%=============================================================================

SLURM Array Job 允许你\textbf{一次提交一批结构相同但参数不同的作业}。核心思路很简单：

\begin{enumerate}[nosep]
  \item 把三维矩阵（环境 $\times$ 架构 $\times$ 种子）展开为一个一维的数字索引 \texttt{0, 1, 2, ..., N-1}
  \item SLURM 为每个索引创建一个独立的 Task，通过环境变量 \texttt{SLURM\_ARRAY\_TASK\_ID} 告诉脚本``你是第几号''
  \item 脚本根据这个编号反推出对应的（环境, 架构, 种子）组合
\end{enumerate}

假设有 $E$ 个环境、$A$ 个架构、每组 $N$ 个种子，总任务数 $= E \times A \times N$。

\begin{lstlisting}[language=slurm,title={\texttt{scripts/run\_matrix.sh}}]
#!/bin/bash
#SBATCH -J rl_matrix
#SBATCH --time=00-12:00:00
#SBATCH -p gpu,preempt
#SBATCH -N 1
#SBATCH -n 1
#SBATCH --cpus-per-task=4
#SBATCH --mem=16g
#SBATCH --gres=gpu:1
#SBATCH --output=logs/%x_%A_%a.out
#SBATCH --error=logs/%x_%A_%a.err
#SBATCH --mail-type=END,FAIL
\end{lstlisting}

注意 \texttt{--output} 里用了 \texttt{\%A}（主 Job ID）和 \texttt{\%a}（Task 序号），这样每个 Task 有独立的日志文件。

\begin{lstlisting}[language=slurm]
# ---- 配置 ----
NUM_SEEDS=5

# 读取环境和架构列表到数组
mapfile -t ENVS < configs/envs.txt
mapfile -t ARCHS < configs/archs.txt

NUM_ENVS=${#ENVS[@]}
NUM_ARCHS=${#ARCHS[@]}
\end{lstlisting}

\texttt{mapfile} 把文本文件逐行读入 Bash 数组。这样我们可以用索引来访问每个元素。

\begin{lstlisting}[language=slurm]
# ---- 从 SLURM_ARRAY_TASK_ID 解码实验参数 ----
# 把一维索引还原为三维坐标，类似于把数组下标转换为行列号
TASK_ID=${SLURM_ARRAY_TASK_ID}

SEED=$((TASK_ID % NUM_SEEDS))          # 最内层：种子
TASK_ID=$((TASK_ID / NUM_SEEDS))
ARCH_IDX=$((TASK_ID % NUM_ARCHS))      # 中间层：架构
ENV_IDX=$((TASK_ID / NUM_ARCHS))       # 最外层：环境

ENV=${ENVS[$ENV_IDX]}
ARCH=${ARCHS[$ARCH_IDX]}

echo "=== Task ${SLURM_ARRAY_TASK_ID}: env=${ENV}, arch=${ARCH}, seed=${SEED} ==="
\end{lstlisting}

这段代码的逻辑和``把一个一维数组下标转成多维下标''完全一样（取余、整除），如果你学过数组的线性化存储就会觉得很熟悉。

\begin{lstlisting}[language=slurm]
# ---- 加载环境并运行 ----
module load cuda/12.9.0
module load miniforge/25.3.0
conda activate rl_env

python train.py \
  --env "$ENV" \
  --arch "$ARCH" \
  --seed "$SEED" \
  --output-dir "results/${ENV}/${ARCH}/seed_${SEED}"
\end{lstlisting}

\subsubsection{提交脚本}

提交脚本负责计算总任务数，然后一条命令提交所有任务：

\begin{lstlisting}[language=bash,title={\texttt{scripts/submit\_all.sh}}]
#!/bin/bash
# 计算总任务数
NUM_ENVS=$(wc -l < configs/envs.txt)
NUM_ARCHS=$(wc -l < configs/archs.txt)
NUM_SEEDS=5

TOTAL=$((NUM_ENVS * NUM_ARCHS * NUM_SEEDS))
MAX_INDEX=$((TOTAL - 1))

echo "Submitting ${TOTAL} tasks (${NUM_ENVS} envs x ${NUM_ARCHS} archs x ${NUM_SEEDS} seeds)"

# 创建日志目录（必须提前存在，否则 SLURM 写日志会失败）
mkdir -p logs results

# 提交 Array Job
sbatch --array=0-${MAX_INDEX} scripts/run_matrix.sh
\end{lstlisting}

以 5 环境 $\times$ 5 架构 $\times$ 5 种子为例，共 125 个任务，一条命令全部提交。

\subsubsection{控制并发数}

集群对每用户有 GPU 数量限制。用 \texttt{\%} 符号限制同时运行的 Array Task 数：

\begin{lstlisting}[language=bash]
# 最多同时跑 10 个任务（因为 gpu 分区每用户限 10 块 GPU）
sbatch --array=0-124%10 scripts/run_matrix.sh
\end{lstlisting}

如果你同时用了 \texttt{preempt} 分区，上限可以到 20。

%=============================================================================
\subsection{方案二：嵌套循环提交}
%=============================================================================

如果不同实验需要不同的资源配置（比如 MuJoCo 环境需要更多内存），可以用脚本循环提交独立作业，对每组实验单独设置参数：

\begin{lstlisting}[language=bash,title={\texttt{scripts/submit\_loop.sh}}]
#!/bin/bash
NUM_SEEDS=5

while IFS= read -r ENV; do
  while IFS= read -r ARCH; do
    for SEED in $(seq 0 $((NUM_SEEDS - 1))); do
      # 根据环境选择资源配置
      if [[ "$ENV" == *"Cheetah"* || "$ENV" == *"Hopper"* || "$ENV" == *"Walker"* ]]; then
        MEM="32g"
        TIME="01-00:00:00"
      else
        MEM="8g"
        TIME="00-04:00:00"
      fi

      sbatch \
        -J "rl_${ENV}_${ARCH}_s${SEED}" \
        --time=${TIME} \
        -p gpu,preempt \
        --gres=gpu:1 \
        --cpus-per-task=4 \
        --mem=${MEM} \
        --output="logs/${ENV}_${ARCH}_s${SEED}.out" \
        --error="logs/${ENV}_${ARCH}_s${SEED}.err" \
        --wrap="module load cuda/12.9.0 && module load miniforge/25.3.0 && \
                conda activate rl_env && \
                python train.py --env ${ENV} --arch ${ARCH} --seed ${SEED} \
                --output-dir results/${ENV}/${ARCH}/seed_${SEED}"
    done
  done < configs/archs.txt
done < configs/envs.txt
\end{lstlisting}

\texttt{sbatch --wrap} 可以直接传入命令字符串作为作业内容，不需要单独写作业脚本文件。适合逻辑简单的场景。
